\movedsub{Choosing among the three}
\label{sec:choosing}

The recommendations above say how to report a faithfulness result. They do not say which evaluator to run, and nothing in this paper supports a single answer to that, since the three differ on properties that matter different amounts in different applications. Table~\ref{tab:choose} collects what we measured into the form that decision needs. Every entry comes from a table earlier in the paper; none of it is new evidence.

Read it against the deployment, not in the abstract. If answers are short and the question is whether a fabricated value gets caught, the deterministic check dominates and the choice of entailment scorer barely matters. If answers are long, no evaluator here catches most fabrications and the gap between the best and worst of them, 24.1\% against 5.1\% on the verbose generator's NQ cell, is the difference between a weak gate and no gate. If the evaluation population cannot be guaranteed free of refusals, AlignScore is the only one of the three that scores refusals below attempted answers for every generator we ran; under the other two, refusals score above real answers for two generators and far below for the third, so pooling them moves the average in a generator-dependent direction. If systems are being compared and the expected difference is small, the relevant property is resolving power, where AlignScore needs 0.012 to 0.037 against 0.021 to 0.060 for the entailment aggregate at the same sample size.

Those criteria do not agree, and we are not proposing a weighting for them. The point of stating them together is that a choice is being made whenever a faithfulness number is reported, and it is usually made by default.

\begin{table*}[t]
\centering\small
\begin{tabular}{lcccc}
\toprule
 & \textbf{NLI-max} & \textbf{Claim-min} & \textbf{AlignScore} & \textbf{Value check} \\
\midrule
Detects a falsified value, terse generator & 57--71\% & 61--71\% & 75--89\% & \textbf{100\%} \\
\quad \dots verbose generator & 5--14\% & 25--34\% & 24--49\% & \textbf{100\%} \\
Sensitivity to answer length (gap above) & 54 pts & \textbf{36 pts} & 46 pts & none \\
Resolving power, minimum detectable effect at 80\% & .021--.060 & .024--.050 & \textbf{.012--.037} & --- \\
Scores a refusal above a real answer? & 2 of 3 gens & 2 of 3 gens & \textbf{never} & --- \\
Error types covered & any & any & any & symbolic only \\
False-positive rate & --- & --- & --- & 0.8--5.2\% \\
\bottomrule
\end{tabular}
\caption{What we measured about each instrument, arranged for choosing one. Detection is against a falsified numeric value at the fixed gate (Table~\ref{tab:2}); the entity probe gives the same ordering (Table~\ref{tab:entity}). MDE is from Appendix~\ref{app:power}. Refusal behaviour is the mean refusal score against the mean attempted-answer score, per generator. The sensitivity row is the detection gap between the two generators, averaged over the two datasets; claim-min is the least length-sensitive of the three and still gives up 36 points. The deterministic value check covers only symbolic content, which is why it supplements an evaluator instead of replacing one. One column of this table is dataset-conditional: the evaluator ordering it summarises holds on the two datasets where retrieval works and inverts where retrieval fails (Section~\ref{sec:qasper}), so these properties are measured on a working retrieval regime and should not be read as intrinsic. Its two numbers are not the same quantity, and only one is a measurement: the 100\% follows by construction, since the probe defines the substituted value as absent from the context and a containment test must therefore flag it, while 0.8--5.2\% is the measured rate at which values in untouched answers are flagged --- chunk indices and list markers if those are counted, and otherwise quantities the answer derives rather than copies.}
\label{tab:choose}
\end{table*}
